\chapter{Complexity — Transducer}
%%% generated by the blueprint pipeline from TCSlib.Complexity.ClassNP.Transducer
%%%%%%%%%%%%%%%%%%%%%%%%%%%%%%%%%%%%%%%%%%%%%%%%%%%%%%%%%%%%%%%%%%%%%%%%%%%%%%%
\section{Overview}

A \emph{one-pass transducer} (a Mealy machine) reads its binary input once from left to
right in a finite control state, emitting at most one bit per input bit. This module
realizes every such transducer as a Turing machine with no work tapes that runs in time
$|x| + 1$, so its string function is polynomial-time computable. It is the machine-side
tool behind several reductions of [AB09, ch.~6], among them stages of the reduction from
circuit satisfiability to 3SAT [AB09, Lem.~6.11].

%%%%%%%%%%%%%%%%%%%%%%%%%%%%%%%%%%%%%%%%%%%%%%%%%%%%%%%%%%%%%%%%%%%%%%%%%%%%%%%
\section{Declarations}

\begin{definition}[String function of a one-pass transducer]
\lean{Complexity.transduce}
\label{Complexity.transduce}
\leanok
Let $\sigma$ be a set of control states, let $\delta : \sigma \times \{0,1\} \to \sigma$ be a
transition function, and let $o : \sigma \times \{0,1\} \to \{0,1\} \cup \{\bot\}$ be an
emission function, where the value $\bot$ means that nothing is emitted. The string
function of the one-pass transducer $(\delta, o)$ started in state $q \in \sigma$ is the map
$w \mapsto \tau_{\delta,o}(q, w)$ on $\{0,1\}^*$ defined recursively by
$\tau_{\delta,o}(q, \varepsilon) = \varepsilon$ and
\[ \tau_{\delta,o}(q, b\,w) = \bar o(q,b)\;\tau_{\delta,o}(\delta(q,b), w) \]
for a bit $b$ and a string $w$, where $\bar o(q,b)$ is the one-bit word $o(q,b)$ when
$o(q,b) \in \{0,1\}$ and the empty word when $o(q,b) = \bot$.
\end{definition}

\begin{theorem}[Transducer output on a concatenation]
\lean{Complexity.transduce_append}
\label{Complexity.transduce_append}
\leanok
\uses{Complexity.transduce}
\difficulty{3}
Let $\sigma$ be a set of states, $\delta : \sigma \times \{0,1\} \to \sigma$ a transition
function and $o : \sigma \times \{0,1\} \to \{0,1\} \cup \{\bot\}$ an emission function, and
write $\tau_{\delta,o}(q, w)$ for the output of the one-pass transducer that reads $w$ from
left to right starting in state $q$, where each bit $b$ read in state $p$ emits the bit
$o(p,b)$ (or nothing if $o(p,b) = \bot$) and moves to state $\delta(p,b)$. Then for every
state $q$ and all bit strings $u, v \in \{0,1\}^*$,
\[ \tau_{\delta,o}(q, uv) = \tau_{\delta,o}(q, u)\;\tau_{\delta,o}\bigl(\hat\delta(q,u), v\bigr), \]
where $\hat\delta(q,u)$ is the state reached from $q$ by applying $\delta$ successively to
the bits of $u$ from left to right.
\end{theorem}

\begin{definition}[Transition function of the transducer machine]
\lean{Complexity.transducerTr}
\label{Complexity.transducerTr}
\leanok
\uses{Turing.Action}
Let $\sigma$ be a set of states, $\delta : \sigma \times \{0,1\} \to \sigma$ a transition
function and $o : \sigma \times \{0,1\} \to \{0,1\} \cup \{\bot\}$ an emission function
($\bot$ meaning no output). This defines the transition function, for a binary Turing
machine with state set $\sigma$, a read-only input tape, no work tapes and a write-only
output tape, given as follows. In state $q$ reading an input bit $b$, the machine appends
$o(q,b)$ to the output tape (nothing if $o(q,b) = \bot$), moves its input head one cell
to the right and enters state $\delta(q,b)$; in any state, reading a blank input cell, it
leaves the head in place, writes nothing and halts.
\end{definition}

\begin{definition}[Turing machine running a one-pass transducer]
\lean{Complexity.transducerTM}
\label{Complexity.transducerTM}
\leanok
\uses{Complexity.transducerTr, Turing.FinTM}
Let $\sigma$ be a finite set of states, $\delta : \sigma \times \{0,1\} \to \sigma$ a
transition function, $o : \sigma \times \{0,1\} \to \{0,1\} \cup \{\bot\}$ an emission
function ($\bot$ meaning no output) and $q_0 \in \sigma$ a start state. The transducer
machine $M_{\delta,o,q_0}$ is the finite binary Turing machine with no work tapes, state
set $\sigma$ and start state $q_0$ whose transition function, in state $q$ reading an input
bit $b$, appends $o(q,b)$ to the output tape (nothing if $\bot$), moves the input head one
cell right and enters state $\delta(q,b)$, and on reading a blank input cell halts without
moving or writing.
\end{definition}

\begin{theorem}[Run of the transducer machine on a suffix]
\lean{Complexity.transducer_run}
\label{Complexity.transducer_run}
\leanok
\uses{Complexity.transduce, Complexity.transducerTM, Complexity.transducerTr, Turing.Action.apply, Turing.Cfg, Turing.Cfg.inputSymbol, Turing.FinTM.inputSymbol_at, Turing.MultiTapeTM.runFrom, Turing.MultiTapeTM.runFrom_succ_eq_step, Turing.MultiTapeTM.runFrom_zero, Turing.MultiTapeTM.step, Turing.moveInputPos, Turing.moveInputPos_pos_of_ne_right}
\difficulty{5}
Let $\sigma$ be a finite set of states, $\delta : \sigma \times \{0,1\} \to \sigma$,
$o : \sigma \times \{0,1\} \to \{0,1\} \cup \{\bot\}$ and $q_0 \in \sigma$, and let $M$ be
the binary Turing machine with no work tapes, state set $\sigma$ and start state $q_0$ that,
in state $p$ reading an input bit $b$, appends $o(p,b)$ to its output (nothing if
$o(p,b) = \bot$), moves its input head one cell right and enters $\delta(p,b)$, and that
halts on reading a blank. Let $x = x_0 x_1 \cdots x_{n-1} \in \{0,1\}^*$, let
$0 \le i \le n$, let $w = x_i x_{i+1} \cdots x_{n-1}$ be the suffix of $x$ starting at
position $i$, and let $C$ be any configuration of $M$ on input $x$ that is in state
$q \in \sigma$ with its input head on the cell holding $x_i$ (on the blank just right of
the input when $i = n$). Then running $M$ for $|w| + 1$ steps from $C$ reaches the halting
state, and the output tape then holds the output of $C$ followed by
$\tau_{\delta,o}(q, w)$, where $\tau_{\delta,o}(q, w)$ denotes the word emitted by reading
$w$ from left to right starting in state $q$, each bit $b$ read in state $p$ emitting
$o(p,b)$ (or nothing) and moving to $\delta(p,b)$.
\end{theorem}

\begin{theorem}[Transducer machine computes the transducer function]
\lean{Complexity.transducerTM_computes}
\label{Complexity.transducerTM_computes}
\leanok
\uses{Complexity.transduce, Complexity.transducerTM, Complexity.transducer_run, Turing.FinTM.ComputesFunInTime, Turing.FinTM.computesInTime_iff, Turing.MultiTapeTM.initCfg}
\difficulty{3}
Let $\sigma$ be a finite set of states, $\delta : \sigma \times \{0,1\} \to \sigma$ a
transition function, $o : \sigma \times \{0,1\} \to \{0,1\} \cup \{\bot\}$ an emission
function ($\bot$ meaning no output) and $q_0 \in \sigma$, and let $M$ be the binary Turing
machine with no work tapes, state set $\sigma$ and start state $q_0$ that, in state $p$
reading an input bit $b$, appends $o(p,b)$ to its output (nothing if $\bot$), moves its
input head right and enters $\delta(p,b)$, and halts on reading a blank. Then $M$ computes
the function $x \mapsto \tau_{\delta,o}(q_0, x)$ in time $n + 1$: on every input
$x \in \{0,1\}^*$ it halts within $|x| + 1$ steps with $\tau_{\delta,o}(q_0, x)$ on its
output tape, where $\tau_{\delta,o}(q_0, x)$ is the word emitted by reading $x$ from left to
right starting in $q_0$, each bit $b$ read in state $p$ emitting $o(p,b)$ (or nothing) and
moving to $\delta(p,b)$.
\end{theorem}

\begin{theorem}[One-pass transducers are polynomial-time computable]
\lean{Complexity.polyTimeComputable_transduce}
\label{Complexity.polyTimeComputable_transduce}
\leanok
\uses{Complexity.PolyTimeComputable, Complexity.polyTimeComputable_of_linear, Complexity.transduce, Complexity.transducerTM, Complexity.transducerTM_computes}
\difficulty{2}
Let $\sigma$ be a finite set of states, $\delta : \sigma \times \{0,1\} \to \sigma$ a
transition function, $o : \sigma \times \{0,1\} \to \{0,1\} \cup \{\bot\}$ an emission
function ($\bot$ meaning no output) and $q_0 \in \sigma$, and let $\tau_{\delta,o}(q_0, x)$
be the word emitted by the one-pass transducer reading $x \in \{0,1\}^*$ from left to right
starting in $q_0$, where each bit $b$ read in state $p$ emits $o(p,b)$ (or nothing) and
moves to $\delta(p,b)$. Then the map $x \mapsto \tau_{\delta,o}(q_0, x)$ is
polynomial-time computable: some finite binary Turing machine computes it within
$C(n+1)^c$ steps on every input of length $n$, for some constants $C, c \in \mathbb{N}$.
\end{theorem}
